\scanpage{86}
\mspage{17}\uline{Formes positives}

Le \unsure{théor.|thème|théorie} étant géométrique, nous nous plaçons sur un corps alg. clos $k$. Nous ne considérons que les motifs $1/2$ \uline{simples} sur $k$, \inserted{purs,} de poids $p$. Soit $M$ un tel motif, considérons les formes
\[
M\times M\longrightarrow\mathbb{Q}(-p)
\]
i.e. les éléments de
\[
\Gamma_k\bigl(\mathscr{H}om(M\otimes M,\mathbb{Q}(-p))\bigr)
=\unsure{\Phi|\mathscr{L}|\mathcal{L}}(M)
\]
où le \deleted{\unsure{$\Gamma_k$ est}} $\mathscr{H}om$ est bien de poids $0$. Elles forment un vectoriel de \deleted{poids \unsure{0}}\inserted{dim finie} sur $\mathbb{Q}$. On considère le \inserted{sous-}espace \deleted{des} \deleted{\unsure{formé}} des formes \uline{symétriques} si $p$ pair, \uline{alternées} si $p$ impair. Cet espace sera noté $\Phi^s(M)$. Noter les hom canoniques
\[
\Phi^s(M)\otimes\Phi^s(N)\unsure{\longrightarrow|\mathbin{-}}\Phi^s(M\otimes N)
\]
% archive sheet number: 85
